\documentclass[pdflatex,sn-apa]{sn-jnl}

%%%% Standard Packages
%%<additional latex packages if required can be included here>

\usepackage{graphicx}%
\usepackage{multirow}%
\usepackage{amsmath,amssymb,amsfonts}%
\usepackage{amsthm}%
\usepackage{mathrsfs}%
\usepackage[title]{appendix}%
\usepackage{xcolor}%
\usepackage{textcomp}%
\usepackage{manyfoot}%
\usepackage{booktabs}%
\usepackage{algorithm}%
\usepackage{algorithmicx}%
\usepackage{algpseudocode}%
\usepackage{listings}%

\raggedbottom
%%\unnumbered% uncomment this for unnumbered level heads

\begin{document}

\title[Article Title]{Crypto Fraud Detection}

\author{\fnm{Can-Elian} \sur{Barth}}

\author{\fnm{Florian} \sur{Baumgartner}}

\author{\fnm{Aaron} \sur{Brülisauer}}

\abstract{
\textbf{Purpose:} The rapid growth of cryptocurrencies has been accompanied by a rise in fraudulent activities, presenting significant challenges to financial systems and regulatory frameworks. This research aims to detect fraudulent activity in cryptocurrency markets using financial data and social media interactions.

\textbf{Methods:} The study utilizes open and close price data, alongside embeddings generated from Reddit and Twitter posts and comments, as primary features for predictive modelling. MultiRocket served as the main machine learning algorithm, effectively identifying patterns indicative of fraudulent behaviour. Attempts were made to implement LSTM models and the ConvTran architecture (via \href{https://github.com/Navidfoumani/ConvTran}{ConvTran}), though these efforts were not completed.

\textbf{Results:} MultiRocket successfully detected potential fraudulent activity based on the provided data features, showcasing its suitability for time series and embedding-based analysis. The unfinished implementation of LSTM and ConvTran models highlighted the complexity of these approaches and the need for further research.

\textbf{Conclusion:} This study lays the groundwork for future advancements in fraud detection for cryptocurrency ecosystems. The use of financial and social media data for machine learning-driven analysis offers a promising direction, and ongoing efforts to refine and expand model capabilities can further enhance detection accuracy and scalability.
}

\keywords{Cryptocurrencies, Bitcoin, Ethereum, Fraud, Web Data Acquisition, Social Media, Reddit, X (Twitter), Time Series Analysis, MultiRocket, LSTM, Deep Learning}

\maketitle

\section{Introduction}\label{intro}

Others did research fraudulent transaction and price manipulation like Rug Pulls \citep{mazorra_do_not_rug_on_me_2022}.

Research topic in context of scientific world

motivation

research questions that this project aims to answer

\subsection{Related Work}\label{related_work}
What did others do?

\subsection{Preliminaries}\label{preliminaries}

What knowledge is useful to read this paper?

\subsubsection{Crypto Exchanges}\label{crypto_stocks}

What are crypto exchanges and what is different from normal exchanges?

\subsubsection{Embedding}\label{embedding}

Why and how were the posts/comments embedded?

\subsubsection{MultiRocket}\label{multirocket}

What is MultiRocket?

\section{Methods}\label{methods}

1. gather data

2. use it

\subsection{Data Collection}\label{data_collection}

What was the data-gathering approach?

\subsubsection{Reddit Posts/Comments}\label{reddit}

What data was gathered?

\subsubsection{X (Twitter) Posts}\label{twitter}

What was scraped from X?

\subsubsection{Price Data}\label{price}

What price data was used and where was it from?

\subsection{Coin Labeling}\label{coin_labeling}

How was the coins labeled as fraud / non-fraud?

\subsection{Scam Detection}\label{scam_detection}

What is the idea, how scams are detected? -> ML

\subsubsection{Data Pipeline}\label{data_pipeline}

What was done after downloading the posts/comments and price data until np-arrays/torch-tensors were ready?

\subsubsection{Machine Learning}\label{machine_learning}

What was done on train data and how did it perform on val coins?

\section{Results}\label{results}

\subsection{Data Gathered / EDA}

What and how much data in which time ranges was gathered?

What insights could be found with EDA?

\subsection{ML-Performance on Test Set}

What was achieved on the test set?

\section{Discussion}\label{discussion}

Discuss the results and what are the conclusions (interpret this findings in context)?

\subsection{Limitations}\label{limitation}

In which ways were this approach and this project limited?

\subsection{Future Work}\label{future_work}

Ideas for the future? How could they be tested/verified?

\subsubsection{LSTM}\label{lstm}

What was done and what would have to be done to finalize it?

\subsubsection{Other Deep Learning Approaches}\label{deep_learning}

What was done and what would have to be done to finalize it?

\backmatter

\section*{Acknowledgements}

We, the authors, would like to extend our special thanks to Gabriel Torres Gamez for his collaboration with us on gathering and conducting exploratory data analysis (EDA) on social media data.

We also wish to express our gratitude to Moritz Kirschmann and Adrian Brändli for their invaluable guidance throughout this Challenge-X. Moritz especially supported us in developing machine learning strategies and guiding the process from raw data to a functional ML model. Adrian provided critical assistance in defining fraud within the cryptocurrency space and addressing the societal embedding of a data science research question within the broader societal context.

\section*{Declarations}

\begin{itemize}
\item AI tools

Parts of the text in this paper have been written or enhanced with the assistance of ChatGPT \citep{OpenAI2024}.

\item Code availability

All code is available on \href{https://github.com/CryptoFraudDetection}{github.com/CryptoFraudDetection}.

\item Author contribution

All authors contributed equally to all aspects of this work.

\end{itemize}

\begin{appendices}

\section{Section title of first appendix}\label{secA1}

An appendix contains supplementary information that is not an essential part of the text itself but which may be helpful in providing a more comprehensive understanding of the research problem or it is information that is too cumbersome to be included in the body of the paper.

\end{appendices}

\bibliography{literature/references,literature/manual_ref}

\end{document}
